\providecommand{\kBKM}{\kappa_{\mathrm{BKM}}}
\providecommand{\kch}{\kappa_{\mathrm{ch}}}
\providecommand{\kcat}{\kappa_{\mathrm{cat}}}
\providecommand{\kfib}{\kappa_{\mathrm{fiber}}}
\providecommand{\Z}{\mathbb{Z}}
\providecommand{\DTred}{Z^{\mathrm{red}}_{\mathrm{DT}}}
\providecommand{\GWred}{Z^{\mathrm{red}}_{\mathrm{GW}}}
\providecommand{\PTred}{Z^{\mathrm{red}}_{\mathrm{PT}}}
\providecommand{\Dfive}{D_5}
\providecommand{\DeltaFive}{\Delta_5}
\providecommand{\PhiTenUn}{\Phi_{10}^{\mathrm{un}}}
\providecommand{\PhiTenOP}{\Phi_{10}^{\mathrm{OP}}}
\providecommand{\chiTen}{\chi_{10}}
\providecommand{\phizo}{\phi_{0,1}}
\providecommand{\gDelta}{\mathfrak{g}_{\Delta_5}}

\section*{OP scalar normalisation and Borcherds denominator for \(K3\times E\)}

Let \(S\) be a projective K3 surface, \(E\) an elliptic curve, and
\(X=S\times E\). Five objects meet at the same Siegel variable:
the primitive \(N=1\) reduced curve-counting theorem of
Oberdieck-Pixton, the Bryan-Oberdieck CHL scalar conjecture, the
Pandharipande-Thomas KKV theorem, the fibre-class theorem of
Oberdieck-Shen, and the Borcherds-Gritsenko-Nikulin denominator
construction. Each has its own hypotheses and output. Their common
automorphic image is the Igusa square attached to the weight-five
denominator \(\DeltaFive\).

\subsection*{Normalisation}

The Eichler-Zagier weak Jacobi form in the Gritsenko-Nikulin
normalisation has
\[
  \phizo(\tau,z)
  =
  (y^{-1}+10+y)+O(q).
\]
The K3 elliptic genus is \(2\phizo\). The Borcherds denominator
\(\DeltaFive\) uses the half-index input \(\phizo\), so the Borcherds
constant is \(c_1(0)=10\), and Borcherds' weight formula gives
\[
  \kBKM(\DeltaFive)=\frac{c_1(0)}{2}=5.
\]
Gritsenko and Nikulin use the monic product
\[
  \Dfive(Z)=64^{-1}\DeltaFive(Z).
\]
The unnormalised Igusa square and the OP scalar square are therefore
\[
  \PhiTenUn=\DeltaFive^2,\qquad
  \PhiTenOP=\Dfive^2=4096^{-1}\DeltaFive^2.
\]
Thus the OP reduced Donaldson-Thomas scalar in the primitive chamber is
\[
  \DTred(X)
  =
  -(\PhiTenOP)^{-1}
  =
  -\Dfive^{-2}
  =
  -4096\,\DeltaFive^{-2}.
\]
The schematic expression \(-C/\DeltaFive^2\) has \(C=4096\) in this
normalisation. The bare expression \(-\DeltaFive^{-2}\) denotes the
unnormalised convention.

\subsection*{Primitive \(N=1\) scalar theorem}

Oberdieck and Pixton prove the primitive reduced Gromov-Witten theory
of \(S\times E\) and deduce the Igusa cusp-form identity. For a
primitive class
\(\beta_h\in H_2(S,\Z)\) with
\(\langle\beta_h,\beta_h\rangle=2h-2\), their partition function is the
Laurent expansion
\[
  \GWred(S\times E;\beta_h)=-\chiTen^{-1},
  \qquad
  \chiTen=\PhiTenOP.
\]
Oberdieck's reduced DT/PT comparison and the
Oberdieck-Pandharipande GW/PT change of variables identify the same
Laurent expansion with the Behrend-weighted reduced DT/PT scalar:
\[
  \DTred(S\times E;\beta_h\ \mathrm{primitive})
  =
  \PTred(S\times E;\beta_h\ \mathrm{primitive})
  =
  -4096\,\DeltaFive^{-2}.
\]
This theorem is scalar: it fixes the Igusa factor, the monic OP
normalisation, and the Behrend sign. Hall algebra,
Drinfeld-double, and BKM-recognition structures require additional
source data.

\subsection*{Primitive CHL weights}

For a symplectic K3 automorphism \(g_N\) of order
\(N\in\{1,2,3,4,6\}\) and an order-\(N\) translation \(t_N\) of \(E\),
the diagonal quotient
\[
  X_N=(S\times E)/\langle(g_N,t_N)\rangle
\]
is a smooth elliptic CHL threefold. The primitive CHL Borcherds
constants are
\[
\begin{array}{c|c|c|c}
N & \text{frame shape} & c_N(0) & \kBKM(\Phi_N)=c_N(0)/2\\
\hline
1 & 1^{24} & 10 & 5\\
2 & 1^8 2^8 & 8 & 4\\
3 & 1^6 3^6 & 6 & 3\\
4 & 1^4 2^2 4^4 & 4 & 2\\
6 & 1^2 2^2 3^2 6^2 & 2 & 1 .
\end{array}
\]
This table is the primitive CHL ladder. The rows \(N=5,7,8\), when
used in the Gritsenko-Clery diagonal-divisor catalogue or in Nikulin
symplectic-order geometry, have different host geometry and belong
outside the primitive CHL ladder.

\subsection*{Bryan-Oberdieck CHL scalar conjecture}

For an elliptic CHL quotient
\[
  X_N=[S\times E/\langle(g_N,t_N)\rangle],
\]
Bryan and Oberdieck define the primitive reduced DT series
\[
  Z^{X_N}(q,t,p)
  =
  \sum_{h,d,n}
  \mathrm{DT}^{X_N}_{n,(\beta_h,d)}
  q^{d-1}t^{\langle\beta_h,\beta_h\rangle/2}(-p)^n .
\]
Their Conjecture A predicts
\[
  Z^{X_N}(q,t,p)=-\widetilde\Phi_N(\tau,z,\sigma)^{-1},
  \qquad
  q=e^{2\pi i\tau},\quad t=e^{2\pi i\sigma},\quad p=e^{2\pi iz},
\]
where \(\widetilde\Phi_N\) is the CHL Siegel form obtained as the
Borcherds lift of the corresponding twisted-twined K3 elliptic genus.
For \(N=1\), \(\widetilde\Phi_1=\chiTen=\PhiTenOP\). For \(N\ge2\),
the full primitive scalar identity remains conjectural. Bryan and
Oberdieck prove base cases by the orbifold topological vertex and by
degeneration to \(K3\times\mathbb P^1\). The imprimitive sectors are
governed by multiple-cover conjectures, except in the special cases
proved by Oberdieck-Shen.

\subsection*{Ancillary theorems}

Pandharipande and Thomas prove the KKV formula for reduced K3 curve
counts in every genus and every K3 curve class:
\[
  \sum_{g,h}
  (-1)^g r_{g,h}
  (y^{1/2}-y^{-1/2})^{2g}q^h
  =
  \prod_{n\ge1}
  \frac{1}{(1-q^n)^{20}(1-yq^n)^2(1-y^{-1}q^n)^2}.
\]
This theorem controls K3 BPS numbers. Passing from K3 to
\(K3\times E\) uses the reduced \(E\)-translation quotient and the
Oberdieck-Pixton scalar normalisation
above.

Oberdieck and Shen compute all \(E\)-reduced DT invariants in fibre
classes \((0,d)\) on \(S\times E\) and prove special imprimitive checks
through motivic Hall integration to the ring of dual numbers. The full
imprimitive \(K3\times E\) multiple-cover formula remains conjectural
outside those checked sectors.

\subsection*{Borcherds and Hall separation}

Write
\[
  \phizo(\tau,z)=\sum_{n,\ell} f(n,\ell)q^n y^\ell .
\]
The Borcherds-Gritsenko-Nikulin side starts from \(\phizo\), constructs
the denominator \(\DeltaFive\), and gives the BKM superalgebra
\(\gDelta\) with root superdimensions
\[
  \operatorname{sdim}\gDelta_{\alpha}=f(nm,\ell),
  \qquad
  \alpha=(n,\ell,m).
\]
The scalar Oberdieck-Pixton theorem uses the square product attached to
the full K3 elliptic genus \(2\phizo\), hence the reciprocal square
\(-\Dfive^{-2}\). The BKM datum remains the square-root denominator
\(\DeltaFive\) with
\[
  \kBKM(\gDelta)=\kBKM(\DeltaFive)=5.
\]

The reduced DT scalar determines a partition function. A Hall algebra
presentation, a Drinfeld double, and compact Hall-Borcherds recognition
for \(Y^+(K3\times E)\) require separate generator, relation, chamber,
and orientation comparisons. The equality of scalar products supplies a
necessary character check for the positive half.

For \(d=3\), \(\Phi_3\) lands in an \(E_1\)-chiral algebra. Braided
\(E_2\) structure belongs to the centre, double, module, or
representation layer. In the same convention,
\(\mathrm{CoHA}(\mathbb C^3)=Y^+\); the
full \(\mathcal W_{1+\infty}\) belongs after double, centre,
completion, or evaluation. The Hall-Drinfeld/BKM branch is distinct
from the Mukai self-mirror Yangian branch.

\subsection*{Invariant and recognition boundary}

The compact \(K3\times E\) invariants are
\[
  \kcat(K3\times E)=0,\qquad
  \kch(K3\times E)=0,\qquad
  \kch^{\mathrm{Heis}}(K3\times E)=3,\qquad
  \kBKM(\DeltaFive)=5,\qquad
  \kfib(K3)=24.
\]
The first equality is
\[
  \chi(\mathcal O_{K3\times E})
  =
  \chi(\mathcal O_{K3})\chi(\mathcal O_E)
  =
  2\cdot 0
  =
  0.
\]
The compact chiral supertrace is the same alternating Hodge column:
\[
  \kch(K3\times E)
  =
  \sum_{q=0}^3 (-1)^q h^{0,q}(K3\times E)
  =
  1-1+1-1
  =
  0.
\]
The Heisenberg-Mukai specialisation supplies the separate
\(\kch^{\mathrm{Heis}}=3\). The Borcherds weight is \(5\). The fibre
rank is \(\operatorname{rk}H^\bullet(K3,\Z)=24\). These are four
independent constructions.

The \(N=1\) scalar theorem is Oberdieck-Pixton with the normalisation
above. The \(N\ge2\) primitive CHL scalar identity is the
Bryan-Oberdieck conjecture with proved base coefficients. K3 BPS
numbers are controlled by Pandharipande-Thomas KKV. Fibre-class and
special imprimitive reduced DT sectors are controlled by
Oberdieck-Shen. Hall-Borcherds recognition is the separate algebraic
comparison: positive-half generators, BKM root multiplicities,
Drinfeld-double completion, orientation character, and chamber
compatibility.
